%---
% publish: true
% tags: [guide, latex]
%---%
%
% This block is Vellum's frontmatter. Both fences are LaTeX COMMENTS, so
% pdflatex never sees it — which is the deal this whole file is built on:
% everything Vellum adds must be invisible to the other toolchain.

\documentclass{article}
\usepackage{vellum}   % get it from your own instance: /api/vellum.sty
\usepackage{amsmath}
\usepackage{graphicx}

\newcommand{\R}{\mathbb{R}}

\title{LaTeX in Vellum}
\author{The vault}

\begin{document}
\maketitle

\begin{abstract}
  A \texttt{.tex} file is a note here, not an import. It sits in the tree, it
  is searchable by its words, it draws edges in the graph, it can be published
  to the blog, and it still compiles with \texttt{pdflatex} exactly as it did
  before you dropped it in.
\end{abstract}

\section{Linking out}\label{sec:linking}
The recommended form is a real macro, so the file stays portable:
\note{Wikilinks \& Backlinks} links by note title, and
\note[the graph view]{Graph View} gives it different display text. A link the
PDF must not show rides in a comment instead:
%% [[Embeds & Transclusion]] %%

Nothing else had to be invented. \texttt{\textbackslash input}, \texttt{\textbackslash cite},
\texttt{\textbackslash ref} and \texttt{\textbackslash eqref} already say what they mean, so Vellum
just extends their search path to the vault --- local definitions first, always,
so an existing project cannot change how it compiles by being moved here.

\section{Maths}\label{sec:maths}
Inline maths reads as $e^{i\pi} + 1 = 0$, and a display equation is numbered
and addressable:
\begin{equation}\label{eq:euler}
  \int_{-\infty}^{\infty} e^{-x^{2}}\,\mathrm{d}x = \sqrt{\pi}.
\end{equation}
Equation~\eqref{eq:euler} can be pointed at from anywhere in the vault --- a
markdown note writes \texttt{[[LaTeX Notes\#eq:euler]]} to link it, or
\texttt{![[LaTeX Notes\#eq:euler]]} to pull the equation itself into its own
page. A heading and a \texttt{\textbackslash label} are the same kind of thing here, so that
works in both directions and in both formats.

\section{The rest}\label{sec:rest}
Lists, emphasis, quotes and figures render as you would expect:
\begin{itemize}
  \item \emph{Emphasis}, \textbf{bold}, \texttt{monospace}, \textsc{small caps}.
  \item Sections are numbered, and the outline panel follows the hierarchy.
  \item $\R$ and any other \texttt{\textbackslash newcommand} you define reach the maths.
\end{itemize}

\begin{theorem}[Honest boundaries]\label{thm:honest}
  Anything this reader does not implement renders as a quiet marker --- never
  as raw source, and never as a crash. The README lists the supported subset.
\end{theorem}

\end{document}
